\documentclass[10pt,a4paper]{amsart}
\usepackage[margin=19mm]{geometry}
\usepackage{amsmath,amssymb,mathtools}
\usepackage[scr=rsfs,cal=euler]{mathalfa}
\usepackage{xfrac}
\usepackage[unicode,hidelinks,bookmarks=false]{hyperref}
\newcommand{\ie}{i.e.\ }
\newcommand{\cf}{cf.\ }
\newcommand{\resp}{respectively\ }
\newcommand{\Subsection}{\subsection}
\providecommand{\nobreakdash}{\nobreak\hbox{-}}
\setlength{\emergencystretch}{2em}
\multlinegap0pt
\usepackage{amssymb}
\usepackage{float}
\newtheorem{theorem}{Theorem}
\title{Disjoint strong transitive operators on Hilbert $C^{\ast}$-modules}
\author{Song-Ung Ri, Hyon-Hui Ju*}
\date{Source: arXiv:2609.02281v1 (CC BY 4.0)}
\begin{document}
\maketitle

\noindent\emph{Source and licence.} Disjoint strong transitive operators on Hilbert $C^{\ast}$-modules. Version arXiv:2609.02281v1 (CC BY 4.0), distributed by arXiv under the Creative Commons Attribution 4.0 International licence, http://creativecommons.org/licenses/by/4.0/, as recorded in the arXiv licence metadata for this exact version and shown on the abstract page https://arxiv.org/abs/2609.02281v1. Full licence notice: \texttt{licenses/arxiv-2609.02281-CC-BY-4.0.txt}.\par\smallskip

\noindent\textit{Editorial scope.} Counted unit: the article's theorem the4.1 (source0.tex lines 363--379). The article's complete original proof of it is delivered with it (source0.tex lines 381--501, one \texttt{proof} environment of 121 lines). No mathematics is written, repaired or re-derived: the statement and the proof are the article's own lines, in the article's own notation, and no retained line is substituted or re-flowed.  No further source-local context is needed: every cross-reference of every retained line resolves inside the two delivered spans.  Nothing was omitted from any retained span, so the package carries no omission record.  No retained line contains a citation, so the wrapper carries no bibliography and no bibliography entry.  No retained line contains a figure environment, an image-inclusion command or drawing code, so no image is delivered and the package declares no asset dependency.  Editorial formatting only, in addition: an AMS \texttt{amsart} preamble that restates the article's own declarations and macros used by the retained lines, namely the environments \texttt{theorem} on separate counters, together with the standard packages those lines need.  The retained environments print in this document as Theorem 1 in order of appearance, whereas the article prints the same items with the numbering of its own full document.  The complete native source of the article as distributed in the arXiv e-print for this exact version is bundled unchanged as \texttt{sources/209188/source0.tex}.
\begin{theorem}\label{the4.1}
Let $b_{1}$, $b_{2}$, . . . $b_{N}$ be invertible elements of $\mathcal{A}_{1}$ and $\Phi$ be an isometric $\ast$-isomorphism of $\mathcal{A}_{1}$. And let $\mathcal{F}$ be a finitely invariant Furstenberg family and $\{r_{k}\}^{N}_{k=1}\subset \mathbb{N}$ satisfy $0<r_{1}<r_{2}<...<r_{N}$. For $1\leq m\leq N$ we denote $T_{\Phi, b_{m}}:=T_{m}$. Then the following are equivalent.

$(1)$ $T^{r_{1}}_{1}$, $T^{r_{2}}_{2}$, ... ,$T^{r_{N}}_{N}$ are disjoint $\mathcal{F}$-transitive.

$(2)$ For every $p_{\alpha}$ there exist sequences $\{q_{n}\}_{n}$ and $\{d^{(1)}_{n}\}_{n}$, $\{d^{(2)}_{n}\}_{n}$, . . . , $\{d^{(N)}_{n}\}_{n}$ in $\mathcal{A}$ satisfying;
for every $\varepsilon>0$ there exists a set $F_{\varepsilon}\in\mathcal{F}$ such that  for any $n\in F_{\varepsilon}$ and

$(i)$ any $1\leq l\leq N$,
\[\|q_{n}-p^{2}_{\alpha}\|<\varepsilon, \|d^{(l)}_{n}-p^{2}_{\alpha}\|<\varepsilon,\]
\[\|\Phi^{-r_{l}n}(b_{l})\Phi^{-r_{l}n+1}(b_{l})\cdots \Phi^{-1}(b_{l})q_{n}\|<\varepsilon,\]
\[\|\Phi^{r_{l}n-1}(b^{-1}_{l})\Phi^{r_{l}n-2}(b^{-1}_{l})\cdots\Phi^{-1}(b^{-1}_{l})b^{-1}_{l}d^{(l)}_{n}\|<\varepsilon,\]
and

$(ii)$ any $1\leq l\neq s \leq N$,
\[\|\Phi^{r_{s}n-r_{l}n}(b_{l})\Phi^{r_{s}n-r_{l}n+1}(b_{l})\cdots\Phi^{r_{s}n-1}(b_{l})\Phi^{r_{s}n-1}(b^{-1}_{s})\Phi^{r_{s}n-2}(b^{-1}_{s})\cdots \Phi(b^{-1}_{s})b^{-1}_{s}d^{(s)}_{n}\|<\varepsilon.\]
\end{theorem}

\begin{proof}
First we show $(i)\Rightarrow(ii)$. Sinc $T^{r_{1}}_{1}$, $T^{r_{2}}_{2}$,..., $T^{r_{N}}_{N}$ are disjoint $\mathcal{F}$-transitive, there exists a sequence $\{a_{n}\}_{n}\subset \mathcal{A}$ such that 
\[(a_{n}, T^{r_{1}n}_{1}(a_{n}),..., T^{r_{N}n}_{N}(a))\stackrel{\mathcal F}{\longrightarrow} (p_{\alpha}, p_{\alpha},..., p_{\alpha}), \textrm{ as } n\rightarrow \infty.\]
In other words, for every $\varepsilon>0$ there exists a set $F_{\varepsilon}\in \mathcal{F}$ such that for any $n\in F_{\varepsilon}$ and any $1\leq l\leq N$, $\|a_{n}-p_{\alpha}\|<\varepsilon$, $\|T^{r_{l}n}_{l}a_{n}-p_{\alpha}\|<\varepsilon$.
Since $\mathcal{F}$ is finitely invariant, we may assume that for any $n\in F_{\varepsilon}$ $n>N_{\alpha}$. 

Now for $1\leq l\leq N$, we define 
\[q_{n}:=a_{n}p_{\alpha}, d^{(l)}_{n}:=b_{l}\Phi(b_{l})\cdots\Phi^{r_{l}n-1}(b_{l})\Phi^{r_{l}n}(a_{n})p_{\alpha}.\]
Then the sequences $\{q_{n}\}_{n}$, $\{d^{(l)}_{n}\}_{n}\subset \mathcal{A}$, since $\mathcal{A}$ is an ideal of $\mathcal{A}_{1}$. For any $n\in F_{\varepsilon}$ and any $1\leq l, s\leq N, l\neq s$,
\[\|q_{n}-p^{2}_{\alpha}\|=\|a_{n}p_{\alpha}-p^{2}_{\alpha}\|\leq \|a_{n}-p_{\alpha}\|<\varepsilon,\]
\begin{eqnarray}\begin{split}\nonumber
\|d^{(l)}_{n}-p^{2}_{\alpha}\|&=\|b_{l}\Phi(b_{l})\cdot\Phi^{r_{l}n-1}(b_{l})\Phi^{r_{l}n}(a_{n})p_{\alpha}-p^{2}_{\alpha}\|
\\
&=\|T^{r_{l}n}_{l}(a_{n})p_{\alpha}-p^{2}_{\alpha}\|\leq \|T^{r_{l}n}_{l}(a_{n})-p_{\alpha}\|
&<\varepsilon,
\end{split}
\end{eqnarray}

\begin{eqnarray}\begin{split}\nonumber
\|\Phi^{-r_{l}n}(b_{l})\Phi^{-r_{l}n+1}(b_{l})\cdots \Phi^{-1}(b_{l})q_{n}\|&=\|\Phi^{-r_{l}n}(b_{l})\Phi^{-r_{l}n+1}(b_{l})\cdots \Phi^{-1}(b_{l})a_{n}p_{\alpha}\|
\\
&=\|\Phi^{-r_{l}n}(b_{l}\Phi(b_{l})\cdots \Phi^{r_{l}n-1}(b_{l})\Phi^{r_{l}n}(a_{n})-p_{\alpha})p_{\alpha}\|
\\
&\leq \|b_{l}\Phi(b_{l})\cdots \Phi^{r_{l}n-1}(b_{l})\Phi^{r_{l}n}(a_{n})-p_{\alpha}\|=\|T^{r_{l}n}_{l}(a_{n})-p_{\alpha}\|
\\
&<\varepsilon,
\end{split}
\end{eqnarray}

\begin{eqnarray}\begin{split}\nonumber
\|\Phi^{r_{l}n-1}(b^{-1}_{l})\Phi^{r_{l}n-2}(b^{-1}_{l})\cdots\Phi(b^{-1}_{l})b^{-1}_{l}d^{(l)}_{n}\|&=\|\Phi^{r_{l}n}(a_{n})p_{\alpha}\|=\|\Phi^{r_{l}n}(a_{n}-p_{\alpha})p_{\alpha}\|
\\
&=\|a_{n}-p_{\alpha}\|<\varepsilon,
\end{split}
\end{eqnarray}
\begin{eqnarray}\begin{split}\nonumber
\|\Phi^{r_{s}n-r_{l}n}(b_{l})&\Phi^{r_{s}n-r_{l}n+1}(b_{l})\cdots\Phi^{r_{s}n-1}(b_{l})\Phi^{r_{s}n-1}(b^{-1}_{s})\Phi^{r_{s}n-2}(b^{-1}_{s})\cdots \Phi(b^{-1}_{s})b^{-1}_{s}d^{(s)}_{n}\|
\\
&=\|\Phi^{r_{s}n-r_{l}n}(b_{l})\Phi^{r_{s}n-r_{l}n+1}(b_{l})\cdots\Phi^{r_{s}n-1}(b_{l})\Phi^{r_{s}n}(a_{n})p_{\alpha}\|
\\
&=\|\Phi^{r_{s}n-r_{l}n}(b_{l}\Phi(b_{l})\cdots\Phi^{r_{l}n-1}(b_{l})\Phi^{r_{l}n}(a_{n})\Phi^{r_{l}n-r_{s}n}(p_{\alpha}))\|
\\
&=\|b_{l}\Phi(b_{l})\cdots \Phi^{r_{l}n-1}(b_{l}) \Phi^{r_{l}n}(a_{n})\Phi^{r_{l}n-r_{s}n}(p_{\alpha})\|
\\
&=\|T^{r_{l}n}_{l}(a_{n})\Phi^{r_{l}n-r_{s}n}(p_{\alpha})\|
\\
&=\|T^{r_{l}n}_{l}(a_{n})\Phi^{r_{l}n-r_{s}n}(p_{\alpha})-p_{\alpha}\Phi^{r_{l}n-r_{s}n}(p_{\alpha})\|
\\
&\leq\|T^{r_{l}n}_{l}(a_{n})-p_{\alpha}\|
\\
&<\varepsilon,
\end{split}
\end{eqnarray}
which concludes $(i)\Rightarrow(ii)$.

Now we show $(ii)\Rightarrow(i)$. Let $O, V_{1}, ..., V_{N}$ be non-empty subsets of $\mathcal{A}$. Since $\{p^{2}_{\alpha}\}_{\alpha}$ is an approximate unit for $\mathcal{A}$, there exist $p_{\alpha}$ and $x\in O$, $y_{1}\in V_{1}$, . . . , $y_{N}\in V_{N}$ such that $x=p^{2}_{\alpha}x$, $y_{1}=p^{2}_{\alpha}y_{1}$, . . . , $y_{N}=p^{2}_{\alpha}y_{N}$. 
Let us denote $\|x\|+\|y_{1}\|+\cdots+\|y_{N}\|=c$.

Define 
\[x_{n}:=q_{n}x+\sum^{N}_{l=1}S^{r_{l}n}_{l}(d^{(l)}_{n}y_{l}).\]
Now it is sufficient to show that 
\[(x_{n}, T^{r_{1}n}_{1}(x_{n}), ..., T^{r_{N}n}_{N}(x_{N}))\stackrel{\mathcal F}{\longrightarrow} (x, y_{1}, ..., y_{N}), \textrm{ as } n\rightarrow \infty.\] 
From condition (ii), there exist sequences of elements of $\mathcal{A}$ $\{q_{n}\}_{n}$ and $\{d^{(1)}_{n}\}_{n}$, $\{d^{(2)}_{n}\}_{n}$, . . . , $\{d^{(N)}_{n}\}_{n}$  satisfying;

for every $\varepsilon>0$ there exists a set $F_{\varepsilon}\in\mathcal{F}$ such that for any $n\in F_{\varepsilon}$ and any $1\leq l\neq s \leq N$,
\[\|q_{n}-p^{2}_{\alpha}\|<\varepsilon/c, \|d^{(l)}_{n}-p^{2}_{\alpha}\|<\varepsilon/c,\]
\[\|\Phi^{-r_{l}n}(b_{l})\Phi^{-r_{l}n+1}(b_{l})\cdots \Phi^{-1}(b_{l})q_{n}\|<\varepsilon/c,\]
\[\|\Phi^{r_{l}n-1}(b^{-1}_{l})\Phi^{r_{l}n-2}(b^{-1}_{l})\cdots\Phi^{-1}(b^{-1}_{l})b^{-1}_{l}d^{(l)}_{n}\|<\varepsilon/c,\]
\[\|\Phi^{r_{s}n-r_{l}n}(b_{l})\Phi^{r_{s}n-r_{l}n+1}(b_{l})\cdots\Phi^{r_{s}n-1}(b_{l})\Phi^{r_{s}n-1}(b^{-1}_{s})\Phi^{r_{s}n-2}(b^{-1}_{s})\cdots \Phi(b^{-1}_{s})b^{-1}_{s}d^{(s)}_{n}\|<\varepsilon/c.\]
Then 
\begin{eqnarray}\begin{split}\nonumber
\|x_{n}-x\|&=\|q_{n}x+\sum^{N}_{l=1}S^{r_{l}n}_{l}(d^{(l)}_{n}y_{l})-x\|
\\
&\leq \|q_{n}x-p^{2}_{\alpha}x\|+\sum^{N}_{l=1}\|S^{r_{l}n}_{l}(d^{(l)}_{n}y_{l})\|
\\
&\leq\|q_{n}-p^{2}_{\alpha}\|\|x\|+\sum^{N}_{l=1}\|S^{r_{l}n}_{l}(d^{(l)}_{n}y_{l})\|
\\
&=\|q_{n}-p^{2}_{\alpha}\|\|x\|+\sum^{N}_{l=1}\|\Phi^{-1} (b^{-1}_{l})\cdots \Phi^{-r_{l}n} (b^{-1}_{l})\Phi^{-r_{l}n} (d^{(l)}_{n}y_{l})\|
\end{split}
\end{eqnarray}

\begin{eqnarray}\begin{split}\nonumber
&=\|q_{n}-p^{2}_{\alpha}\|\|x\|+\sum^{N}_{l=1}\|\Phi^{r_{l}n}( \Phi^{r_{l}n-1}(b^{-1}_{l})\cdots \Phi(b^{-1}_{l})b^{-1}_{l}d^{(l)}_{n}y_{l})\|
\\
&\leq\|q_{n}-p^{2}_{\alpha}\|\|x\|+\sum^{N}_{l=1}\| \Phi^{r_{l}n-1}(b^{-1}_{l})\cdots \Phi(b^{-1}_{l})b^{-1}_{l}d^{(l)}_{n}\|\|y_{l}\|
\\
&<\frac{\varepsilon}{c}\|x\|+\sum^{N}_{l=1}\frac{\varepsilon}{c}\|y_{l}\|
\\
&=\varepsilon.
\end{split}
\end{eqnarray}
And
\begin{eqnarray}\begin{split}\nonumber
&\|T^{r_{l}n}_{l}(x_{n})-y_{l}\|
=\|T^{r_{l}n}_{l}(q_{n}x)+ T^{r_{l}n}_{l}(\sum^{N}_{s=1}S^{r_{s}N}_{s}(d^{(s)}_{n}y_{s}))-y_{l}\|
\\
&\leq \|T^{r_{l}n}_{l}(q_{n}x)\|+\sum^{N}_{s=1, s\neq l}T^{r_{l}n}_{l}(S^{r_{s}N}_{s}(d^{(s)}_{n}y_{s}))
\\
&=\|\Phi^{r_{l}n}(\Phi^{-r_{l}n}(b_{l})\Phi^{-r_{l}n+1}(b_{l})\cdots \Phi^{-1}(b_{l})q_{n}x)\
\\
&+\sum^{N}_{s=1, s\neq l}\|b_{l}\Phi(b_{l})\cdots\Phi^{r_{l}n-1}(b_{l}) \Phi^{r_{l}n}(\Phi^{-1}(b^{-1}_{s})\cdots\Phi^{-r_{s}n}(b^{-1}_{s})\Phi^{-r_{s}n}(d^{(s)}_{n}y_{s}))\|
\\
&=\|\Phi^{-r_{l}n}(b_{l})\Phi^{-r_{l}n+1}(b_{l})\cdots \Phi^{-1}(b_{l})q_{n}x \|
\\
&+\sum^{N}_{s=1, s\neq l}\|\Phi^{r_{l}n-r_{s}n}(\Phi^{r_{s}n-r_{l}n}(b_{l})\Phi^{r_{s}n-r_{l}n+1}(b_{l})\cdots\Phi^{r_{s}n-1}(b_{l})\Phi^{r_{s}n-1}(b^{-1}_{s})\Phi^{r_{s}n-2}(b^{-1}_{s})\cdots \Phi(b^{-1}_{s})b^{-1}_{s}d^{(s)}_{n}y_{s})\|
\\
&=\|\Phi^{-r_{l}n}(b_{l})\Phi^{-r_{l}n+1}(b_{l})\cdots \Phi^{-1}(b_{l})q_{n}x \|
\\
&+\sum^{N}_{s=1, s\neq l}\|\Phi^{r_{s}n-r_{l}n}(b_{l})\Phi^{r_{s}n-r_{l}n+1}(b_{l})\cdots\Phi^{r_{s}n-1}(b_{l})\Phi^{r_{s}n-1}(b^{-1}_{s})\Phi^{r_{s}n-2}(b^{-1}_{s})\cdots \Phi(b^{-1}_{s})b^{-1}_{s}d^{(s)}_{n}y_{s}\|
\\
&=\|\Phi^{-r_{l}n}(b_{l})\Phi^{-r_{l}n+1}(b_{l})\cdots \Phi^{-1}(b_{l})q_{n}x \|
\\
&+\sum^{N}_{s=1, s\neq l}\|\Phi^{r_{s}n-r_{l}n}(b_{l})\Phi^{r_{s}n-r_{l}n+1}(b_{l})\cdots\Phi^{r_{s}n-1}(b_{l})\Phi^{r_{s}n-1}(b^{-1}_{s})\Phi^{r_{s}n-2}(b^{-1}_{s})\cdots \Phi(b^{-1}_{s})b^{-1}_{s}d^{(s)}_{n}\|\|y_{s}\|
\\
&<\frac{\varepsilon}{c}\|x\|+\sum^{N}_{s=1, s\neq l}\frac{\varepsilon}{c}\|y_{s}\|
\\
&\leq \varepsilon.
\end{split}
\end{eqnarray}
This completes the proof.
\end{proof}

\end{document}
